\section{Transfer Functions}

\subsection{Fundamentals}

Transfer functions map movements in the control space to movements in the display space. In order for this relationship to be meaningful, it must be hardware-independent: all movements must be expressed using standard length and time units such as meters and seconds, not device-specific ones such as Pixels~\cite{casiez2011no}. In the simplest case of a linear relationship, the scale factor between the two spaces is referred to as the Control-Display (CD) gain:

\begin{equation}
    \text{CD gain} = \frac{V_{\text{display}}}{V_{\text{control}}}
\end{equation}

For this coefficient to be unit less, numerator and denominator must be expressed using the same length and time units~\cite{casiez2011no}. Whether static or dynamic, CD gain settings involve a trade-off between gross and fine positioning. High gains reduce the time required to approach a distant target but make it difficult to precisely position the cursor on it. Conversely, low gains support precise positioning but increase the time needed to cover large distances~\cite{mackenzie1995input}. This trade-off is formally captured by Fitts' Law, which predicts the movement time $MT$ required to acquire a target as a function of its distance $D$ and its width $W$, $MT = a + b\,\log_2\!\left(\tfrac{D}{W} + 1\right)$, where the logarithmic term denotes the index of difficulty and $a$ and $b$ are empirically fitted constants. The CD gain can be kept constant or dynamically adjusted over time, for example, based on the kinematics of the control space, i.e.\ as a function of the current input movement velocity~\cite{casiez2011no}. Conceptually, transfer functions can be decomposed into three transformations~\cite{quinn2012exposing, quinn2013touch}. First, translation converts device units such as degrees of rotation, millimeters of finger displacement, or force values into units meaningful to the system, such as pixels or lines. Users may be able to adjust this translation through device controls or system preferences. Second, a gain function amplifies or attenuates the input signal for example, to allow slow precise control when the device is manipulated slowly, and accelerated output when it is manipulated more aggressively. The gain level is typically user-configurable. Third, a persistence component preserves a history of prior inputs and computed parameters such as velocity and acceleration, enabling time-dependent effects such as cumulative gain increases, inertia, or simulated friction. Data from this component can feed back into both translation and gain. Not all of these transformations need to be present in any given transfer function some may be absent, combined, or applied multiple times at different levels of the system stack, for instance when gain is applied both by a driver and again by a UI toolkit~\cite{quinn2012exposing}. This model applies equally to mouse-based pointing and scrolling as well as to touch-based scrolling~\cite{quinn2012exposing, quinn2013touch}. Despite their fundamental role in interaction, transfer functions are typically embedded as black boxes in proprietary operating system drivers and are not publicly documented~\cite{quinn2012exposing, quinn2013touch}. This creates three recurring problems for researchers: first, a lack of knowledge about the current state of the art; second, difficulty in replicating studies, since exact function parameters are unknown or cannot be reproduced and third, the risk of inadvertently introducing confounds into experiments through unknown interactions between the system transfer function and the experimental condition~\cite{quinn2012exposing}. Casiez and Roussel characterise this situation as a practice of bricolage. An unreflective patching together of conditions without precise knowledge of the transfer functions involved~\cite{casiez2011no}. Quinn et al.\ subsequently observed the same problems for scroll-wheel transfer functions~\cite{quinn2012exposing} and for touch scrolling transfer functions on smartphones and tablets~\cite{quinn2013touch}.

\subsection{Pointing Transfer Functions}

\subsubsection{Characterising and Replicating Commercial Functions}

Indirect pointing has been the dominant interaction paradigm on desktop systems for decades. The mouse established itself as the primary input device after outperforming all other devices in early comparative studies on pointing speed and accuracy~\cite{casiez2011no}. In indirect pointing configurations, movements in the control space are mapped to movements in the display space. This mapping is referred to as a transfer function and must be defined in a hardware-independent manner. In the simplest case of a linear mapping, the control-display (CD) gain denotes the scale factor between the speed in the control space and the speed in the display space~\cite{casiez2011no}. Based on the hybrid optimized impulse motor control model, modern transfer functions typically apply high CD gains during the ballistic phase of a pointing movement, corresponding to high input velocities, and low CD gains during the corrective phase, corresponding to low input velocities~\cite{eddy2026tftune}. Despite their widespread use, the concrete design of transfer functions in commercial operating systems has long been poorly understood. \citeauthor{casiez2011no} found that many authors in the literature use constant CD gains as a baseline condition but fail to describe them in sufficient detail to allow replication~\cite{casiez2011no}. Moreover, enforcing a hardware-independent transfer function in practice even a constant one is considerably challenging~\cite{casiez2011no}. To address this, \citeauthor{casiez2011no} developed EchoMouse, an electronic device for characterizing the transfer functions of any system, and libpointing, a toolkit for replicating and comparing the transfer functions used by Windows, OS~X, and Xorg~\cite{casiez2011no}. Using these tools, they showed that the transfer functions of all three systems are non-linear and cannot be approximated by simple low-order polynomials the OS~X functions in particular exhibit piecewise-defined curves, some featuring unexplainable singularities~\cite{casiez2011no}. In a controlled experiment, the default transfer functions of all three systems outperformed a constant CD gain by up to 24\,\% in movement time, with significant differences between systems depending on target size~\cite{casiez2011no}. Notably, none of the participants was subjectively able to perceive any difference between the three system functions~\cite{casiez2011no}.

\subsubsection{Hardware Independence}

A limitation shared by all of these system functions is that they are conventionally expressed in virtual units---pixels per input count---rather than in physical units of distance and speed. \citeauthor{hanada2021relevance} point out that the ``same'' transfer function defined in this way can produce markedly different observable pointer behavior when used with input or output devices of different resolutions: commercial mice range from roughly 400 to 16{,}000 counts per inch, and displays from about 70 to more than 200 dots per inch, so identical motor movements are sensed and rendered at very different scales~\cite{hanada2021relevance}. In a first controlled experiment they show that even small hardware differences measurably affect pointing performance when the function is defined in virtual units, which turns the combination of transfer function and hardware into a potential confound for any experiment that fails to report both precisely~\cite{hanada2021relevance}. As a remedy, they demonstrate the applicability of hardware-independent transfer functions defined in physical units and propose two methods to preserve consistent behavior on operating systems that require hardware-dependent definitions: scaling the function to the resolutions of the connected input and output devices, or selecting the operating-system acceleration setting whose visual behavior most closely matches the intended one~\cite{hanada2021relevance}. A second experiment confirms that both methods yield equivalent pointing performance across different screen and mouse resolutions, and the authors provide a tool to compute equivalent functions between hardware setups so that users can match pointer behavior and researchers can replicate conditions~\cite{hanada2021relevance}.

The practical importance of these software-level factors is underlined by work on touchpads, which have received far less attention than mice despite being the de facto standard input device on laptops. \citeauthor{an2025libpad} note that touchpad user experience is frequently reported to differ across operating systems---with Apple's trackpad widely regarded as superior---yet the underlying reasons have remained scientifically unexamined, as prior work focused mainly on hardware factors such as surface friction and physical size~\cite{an2025libpad}. To isolate the software contribution, they present LibPad, a toolkit that applies OS-specific transfer functions to arbitrary hardware, and reverse engineer the undocumented macOS touchpad function using a private multi-touch framework together with a Fitts'-law pointing task~\cite{an2025libpad}. Unlike mice, which report relative motion counts, touchpads report absolute coordinates and must therefore map input within a bounded surface area, which further couples the function to the device geometry~\cite{an2025libpad}. Their measurements reveal significant performance differences both across different laptop hardware running an identical transfer function and within the same hardware under different function configurations, demonstrating that the transfer function---and not only the physical touchpad---shapes the perceived quality of touchpad pointing~\cite{an2025libpad}.

Because the effective transfer function experienced by a user is thus the product of device firmware, driver, and operating-system settings, achieving a consistent mapping across environments is difficult. \citeauthor{kim2025hardware} address this by moving the transfer function into the device firmware itself: their hardware-embedded approach lets users define a desired function in physical units on the device and then cancels out the influence of the OS-native function and of hardware setting perturbations, so that the uploaded function persists regardless of the operating system or hardware on which the device is used~\cite{kim2025hardware}. Motivated in part by competitive gaming, where consistent aiming behavior across machines is critical, their technical evaluation across transfer functions of various shapes shows robustness and accuracy comparable to the conventional, operating-system-side approach while removing the dependence on external configuration~\cite{kim2025hardware}.

\subsubsection{Automatic and Personalised Optimisation}

Although velocity-dependent transfer functions are known to outperform constant gains, their design has historically relied on manual trial and error or heuristic iteration~\cite{lee2020autogain}. AutoGain~\cite{lee2020autogain} introduced the first method to automatically optimize a gain function from a user's actual motion data, based on the assumption that the aiming error of individual submovements serves as an indicator of the inadequacy of the current gain function: if a submovement undershoots its intended aim point, the gain is too low and should be increased; if it overshoots, the gain is too high and should be decreased~\cite{lee2020autogain}. AutoGain segments each pointing movement into submovements based on local extrema in the speed profile and updates the gain function after each target acquisition trial proportionally to the measured aiming error for the respective speed ranges used~\cite{lee2020autogain}. In a study using a MacBook trackpad, AutoGain produced a gain function after approximately 30 minutes of active use that was comparable in performance to the manually tuned macOS function, a function that had been refined since 1994~\cite{lee2020autogain}. However, AutoGain was not able to significantly outperform existing OS transfer functions~\cite{eddy2026tftune}. TFTune addresses this limitation through a reinforcement learning based approach for automatically personalizing transfer functions. The tuning process is formulated as a Markov Decision Process in which an agent selects a gain value based on the current input speed and receives a reward signal based on the achieved throughput~\cite{eddy2026tftune}. After training, the optimized policy is exported as a discrete lookup table and deployed as the final transfer function~\cite{eddy2026tftune}. Across three evaluation studies, TFTune demonstrated the ability to both create transfer functions from scratch and improve upon pre-established ones: on a MacBook trackpad, it outperformed the macOS default by 7\,\% in movement time within 7 minutes of tuning, and on participants' own Windows devices with varying hardware, an 8\,\% improvement was achieved in as little as one
minute~\cite{eddy2026tftune}. After two to four weeks of real-world daily use, no perceived degradation was reported and all participants chose to retain their personalized functions~\cite{eddy2026tftune}. Unlike AutoGain, which adjusts individual speed bins locally, TFTune optimizes the transfer function as a parameterized policy, yielding smoother and more monotonic functions that generalize better across unseen input speeds~\cite{eddy2026tftune}.

\subsection{Scrolling Transfer Functions}

Although touchscreens had already been gaining traction at the time, \citeauthor{casiez2011no} anticipated that indirect pointing would remain the prevailing paradigm for the foreseeable future. This assumption has since been challenged by the rapid proliferation of mobile devices. As a result, touch-based scrolling, rather than mouse or trackpad interaction, has become the dominant mode of content navigation for a large share of users, making the design and understanding of touch scrolling transfer functions an increasingly critical concern. 

\subsubsection{Input Parameters and the Scrolling Design Space}

A scrolling transfer function can attend to a range of input parameters derived from how users physically interact with the device. These can be classified into three types. Instantaneous measures capture physical properties sensed at a given moment or their derivatives, such as angular velocity or acceleration of a scroll wheel, which can be used to increase scroll magnitude or switch between scrolling modes. Action duration refers to how long an input is sustained or absent, for instance, if scrolling velocity is maintained above a threshold, gain might increase on the assumption that the user intends to travel a long distance. Cumulative and relative inputs introduce memory of prior actions, such as cumulatively adding gain across rapidly repeated wheel rotations in the same direction, canceled if the user pauses or reverses direction~\cite{quinn2012exposing}.

The design space available to a scrolling transfer function is closely tied to the physical characteristics of the input device and the limited mechanics of scroll control. \citeauthor{quinn2012exposing} organise these characteristics by drawing on established input-device taxonomies that classify controls by whether they sense position, motion, or force, along how many linear or rotational dimensions, and whether they operate in absolute or relative terms~\cite{quinn2012exposing}. This framing highlights why scrolling is comparatively impoverished as an input channel: whereas a modern mouse registers thousands of counts per inch across a two-dimensional area that can be traversed without clutching, a typical scroll wheel offers only five or six detents across roughly forty degrees of one-dimensional rotation between clutching actions, and the muscle groups involved impose asymmetric control capabilities across scroll directions~\cite{quinn2012exposing}. Because of this paucity of expressiveness, scrolling functions are likely to attend to several input parameters at once---the rate of device movement, the time between clutched repetitions, and the scroll direction among them---rather than to velocity alone as in pointing~\cite{quinn2012exposing}. To determine which of these parameters commercial systems actually use, and how, \citeauthor{quinn2012exposing} extend the EchoMouse and libpointing approach so that simulated input can be injected at controlled speeds and the resulting system response inspected, turning otherwise opaque driver behavior into measurable transfer functions~\cite{quinn2012exposing}. 

\subsubsection{Reverse Engineering Commercial Scrolling Functions}

Despite this rich design space, existing scrolling transfer functions are typically black boxes embedded in proprietary OS drivers, making them difficult to examine, report, or replicate~\cite{quinn2012exposing}. \citeauthor{quinn2012exposing} adapted the EchoMouse approach of Casiez and Roussel~\cite{casiez2011no} to reverse engineer the scrolling transfer functions of Apple Mac OS~X, Microsoft Windows~7, Microsoft IntelliPoint, and Logitech drivers~\cite{quinn2012exposing}. Their analysis revealed substantial variation across drivers. Neither Windows~7 nor Logitech's SetPoint provide any scroll acceleration (constant gain only), while Mac OS~X applies velocity-dependent gain using a smoothing window of the last eight events to avoid rapid gain changes, and IntelliPoint cumulatively adds gain across successive clutch actions~\cite{quinn2012exposing}. Notably, Mac OS~X exposes a Scrolling Speed slider to users via eight intervals mapped to API values from 0 to 1.7, where a negative value completely disables acceleration for generic devices, or engages a page-scrolling mode for Apple trackpads~\cite{quinn2012exposing}. While scroll-wheel transfer functions can be examined through HID-based methods, touch scrolling presents additional complexity~\cite{quinn2013touch}. While Android's scrolling code is publicly available as open source, Apple and Microsoft require indirect approaches, either programmatic control of the iOS simulator with subsequent replay on an actual device, or decompilation of bytecode in Microsoft's case~\cite{quinn2013touch}.

Regarding the transfer functions themselves, both iOS and Android maintain a 1:1 mapping between finger movement and scrolling output while the finger remains in contact with the display~\cite{quinn2013touch}. Once the finger leaves the surface, however, both systems can substantially amplify the output, applying cumulative gain to assist users in traversing long distances~\cite{quinn2013touch}. The two systems differ in how this gain is applied: iOS applies a capped velocity multiplier after the fourth successive flick gesture (with the cap reaching 16 from the tenth gesture onward), while Android continuously adds any residual scroll velocity to the next flick without an upper cap~\cite{quinn2013touch}. Both systems also apply simulated friction to progressively retard scroll velocity after the finger lifts, with iOS using an exponential decay model and Android relying on spline-based deceleration curves~\cite{quinn2013touch}. These findings highlight several opportunities for improving touch scrolling transfer functions. Since clutch time is a reliable signal of scroll intent, it could be explicitly incorporated as a parameter in gain functions to enable accelerated scrolling across repeated gestures. Furthermore, document-length-dependent gain, which uses a 1:1 mapping at slow input rates but scales gain according to document length at high rates has been identified as a particularly promising approach for mobile devices with small viewports and long content~\cite{quinn2013touch}.

Reverse engineering touch scrolling required different tools than those used for scroll wheels, because touch input cannot be injected through the published USB standards that HID-based methods rely on~\cite{quinn2013touch}. \citeauthor{quinn2013touch} therefore combined controlled software simulation with a high-precision Selective Compliant Articulated Robot Arm (SCARA) that mechanically reproduces touch scrolling gestures with repeatable length and velocity, allowing the black-box functions of commercial devices to be probed directly on the hardware~\cite{quinn2013touch}. Alongside this system-side analysis, they conducted a human-factors experiment to characterise how users actually express scrolling intentions through touch gestures, establishing the practical limits and variability of gesture velocity, length, and repetition that a transfer function must accommodate~\cite{quinn2013touch}. This pairing of a human-behavior study with a device-measurement method is important because it grounds the observed system behavior in the range of inputs users are able and likely to produce, rather than in idealised gestures~\cite{quinn2013touch}. Applying the method to Apple iOS on both iPhone and iPad and to Google Android, they found broad similarities but also substantial platform-specific differences, implying that any experiment involving gestural scrolling is liable to be influenced by the platform on which it is conducted; to support replication and iterative improvement, they released the extracted data tables and the software used to emulate and test the functions~\cite{quinn2013touch}.



\subsubsection{Modelling and Evaluating Scrolling Performance}

A foundational step toward rigorous evaluation of scrolling techniques was made by \citeauthor{hinckley2002quantitative}, who proposed a formal experimental paradigm based on Fitt's Law. While Fitts' Law had traditionally been applied to pointing tasks involving rapid, aimed movement to visible targets, their work demonstrated that it also models scrolling performance well, despite the scrolling target lying beyond the edges of the screen. Their paradigm systematically varies two parameters, scrolling distance and required tolerance. Applying this paradigm to two representative devices, the IBM ScrollPoint isometric joystick and the Microsoft IntelliMouse wheel, revealed a crossover effect in performance as a function of distance: the wheel was faster for short scrolling distances of up to roughly one hundred lines, whereas the isometric joystick became faster for longer distances~\cite{hinckley2002quantitative}. This result is methodologically significant, because a study that reports only overall means---without controlling for distance and tolerance---can reach opposite conclusions depending on the particular distances sampled~\cite{hinckley2002quantitative}. \citeauthor{hinckley2002quantitative} further showed that adding an acceleration algorithm substantially improved the wheel: at a high resolution of one line per notch it matched or exceeded the joystick up to about two hundred lines, and at three lines per notch up to about four hundred lines, demonstrating that scroll acceleration can extend the effective range of an otherwise short-range device~\cite{hinckley2002quantitative}. To date, no ISO standard equivalent to ISO~9241-9 exists for scrolling. 

\citeauthor{zhao2014model} argue that scrolling on touch screens lacked a sufficiently solid theoretical basis, which made it difficult to compare techniques, model user performance, and predict interaction outcomes \cite{zhao2014model}. To address this gap, they do not treat scrolling as a single monolithic action, but first decompose it through a pilot study and task analysis into a simplified target-acquisition problem \cite{zhao2014model}. The authors model 1D scrolling as a two-phase process, consisting of an iterative searching phase followed by a conventional pointing phase once the target becomes visible \cite{zhao2014model}. This decomposition is important because it captures the fact that off-screen target acquisition on touch devices requires an unknown number of repeated pan or flick actions, unlike earlier one- or two-stage pointing paradigms that can largely be described with Fitts' law alone \cite{zhao2014model}. In other words, the study is designed to isolate the repeated clutch-like behavior that characterizes real touch scrolling, rather than collapsing it into a single pointing movement \cite{zhao2014model}. The experiment itself was intentionally constrained to a horizontal 1D setting, with the display window, workspace, target distance, display size, and target width manipulated systematically \cite{zhao2014model}. The authors further separated the interaction into two scrolling modes, panning and flicking, and two feedback conditions, with and without distance feedback, yielding four techniques in total \cite{zhao2014model}. This design choice allowed them to study not only whether scrolling performance differs across techniques, but also how target information and control mode shape the structure of the task \cite{zhao2014model}. Methodologically, the pilot study served to justify the simplifying assumptions used in the model: repeated clutches during search, similar finger-displacement distributions across iterations, a ceiling effect on displacement, and an unknown starting position for the final pointing step \cite{zhao2014model}. The authors then formalize these observations into a quantitative model and validate it experimentally, showing that their task decomposition is not merely descriptive but directly supports the mathematical structure of the model \cite{zhao2014model}. Importantly, they also note that their experimental conditions represent useful extremes, while practical interfaces likely lie somewhere between these boundary cases \cite{zhao2014model}.

From the four assumptions, \citeauthor{zhao2014model} derive a closed-form movement-time model whose central components correspond to the iterative search and the final pointing phase, and they validate both the complete model and its individual components against the experimental data~\cite{zhao2014model}. Beyond the conventional coefficient of determination $R^2$, they compare their model with existing models of related tasks using the Akaike Information Criterion, which balances goodness of fit against model complexity, and report that their formulation performs best under both criteria~\cite{zhao2014model}. The analysis also isolates the effects of the individual factors: flicking and panning structure the search phase differently, and providing distance feedback changes how users regulate their gestures, which in turn affects the number and size of the repeated clutching actions~\cite{zhao2014model}. The authors situate scrolling within the broader significance of the task, citing earlier observations that users may spend on the order of forty minutes merely scrolling during a multi-hour browsing session, so that even small efficiency gains accumulate into meaningful benefits~\cite{zhao2014model}. As a design implication, the model provides a principled basis for comparing and predicting the performance of direct-touch scrolling techniques and offers a foundation on which more sophisticated scrolling interactions can be modelled quantitatively~\cite{zhao2014model}.

\subsubsection{Synthesis and Research Gap}

Taken together, this body of work reveals a consistent trajectory. Early efforts established that transfer functions are pervasive yet poorly documented and introduced tools to characterise and replicate them~\cite{casiez2011no, quinn2012exposing, quinn2013touch}; subsequent work quantified how strongly the function, the hardware, and the operating system jointly determine performance~\cite{hanada2021relevance, an2025libpad, kim2025hardware}; and the most recent contributions turn from characterisation to optimisation, using automated and learning-based methods to tune functions to individual users~\cite{lee2020autogain, eddy2026tftune}. For touch scrolling in particular, however, the transfer functions of commercial platforms remain hand-tuned black boxes, and no comparable method exists for automatically adapting them to individual users or content---the gap that this thesis addresses.
